\scanpage{49}

% B01 | top, continuation of direct limit
et lorsque $i>n$, les \deleted{\unsure{de cette|derniers|deux}} $H^i(\mathfrak g,U(m))$ sont nuls d'après $b)$, \inserted{(car $U(m)$ est fini)} d'où résulte notre assertion.

% B02 | upper quarter, unipotent group
Si enfin $G$ est \deleted{le groupe} un groupe linéaire unipotent connexe, on distingue deux cas

% B03 | upper middle, characteristic zero
1°) Caractéristique nulle. Alors $G$ est $\simeq\bar k^n$ sur $\bar k$, donc aussi sur $k$ (en vertu de $H^1(k,\operatorname{Gl}(\unsure{n|m}))=0$). Et comme $H^i(k,G_a)=0$ pour $i>0$, l'assertion est valable.

% B04 | middle, nonzero characteristic
2°) Car. $p\ne0$. Alors on filtre $G$ par $G\supset pG\supset\ldots\supset p^\ell G=0$. On est ramené au cas où $pG=0$ (les groupes $p^iG$ étant \uline{définis sur $k$}).

% B05 | left, opposite pG = 0
\margin{$//$}

% B06 | middle, descent
Mais on sait qu'alors $G$ est $\simeq\bar k^{\unsure{m|n}}$ sur $\bar k$ Alors \deleted{\unsure{=|\simeq}} aussi sur $k$ (car $H^1(k,\operatorname{Gl}(\unsure{m|n}))=0$), et on finit comme dans 1°.

% B07 | last line, followed by blank half-page
Supposons que $k$ soit quelconque .....

% archive sheet number: 48

